\section{Arithmetic state and determinant comparisons}
The finite boundary point is invariant under the changes proved above.
It is an element of a torsor, with an affine evaluation pairing, and
gives the canonical associated-line vector $[\lambda_S(u),1]$.
A prepared quantum state would additionally require a specified space of
observables and a functional on it. The arithmetic cohomology, the
positive Weil sectors and the analytic determinants therefore leave
precise comparison problems.

\begin{openproblem}[From the classical invariant to a chiral determinant]
\label{problem:acs}
A comparison with Section~\ref{sec:arithmetic-comparison} and
Theorem~\ref{thm:deninger} requires the following constructions.

\begin{enumerate}
\item Fix one cohomological framework: \(X_{\mathrm{et}}\) under the totally
imaginary hypotheses of Theorem~\ref{thm:kim}; ordinary small-site
cohomology with odd-order coefficients when real places are present;
explicitly defined modified cohomology with Tate-cohomological real terms;
or ordinary derived cohomology on \(\widetilde X_{\mathrm{et}}\) with a
specified constructible torsion sheaf, an extension to the site, and
infinite-place specialization maps.  Specify separately any ramification
set and its boundary torsors.

\item For a continuous or non-finite gauge object, define an arithmetic
field groupoid, an action functional, measures or finite sums, state spaces,
gluing maps, and anomaly data.  Assemble them into a symmetric monoidal
functor on a specified arithmetic bordism category.  Kim's finite classical
function does not provide these objects.

\item Construct complexes \(C_{\mathrm{Ar}}\), \(C_{\mathrm{Den}}\), and
\(C_{\mathrm{ACS}}\) in one named derived category and chain maps
\[
 C_{\mathrm{Ar}}\longrightarrow C_{\mathrm{Den}}
 \longrightarrow C_{\mathrm{ACS}}
\]
with stated degrees, completions, quasi-isomorphism hypotheses, and
intertwining relations for the relevant operators.  The first complex must
in particular solve Open Problem~\ref{problem:bar}, and the middle complex
must supply the data assumed in Theorem~\ref{thm:deninger}.

\item To compare a partition object with \(\xi_{\mathbb R}(s)\), construct
the complex parameter \(s\), a holomorphic or meromorphic family of
partition objects, and a proof of the determinant identity.  A single
closed-theory partition number cannot equal a nonconstant function of an
undeclared parameter.

\item A Wilson surface requires higher gauge data: at minimum a crossed
module or \(2\)-group, a \(2\)-connection or arithmetic replacement, and a
surface-holonomy/transgression construction.  It is not defined by an
ordinary finite \(A\)-torsor or by Theorem~\ref{thm:kim}.

\item Any extension to a completed motivic \(L\)-function must begin with
a specified motive, its local factors and archimedean factor, its weight,
and a no-cancellation determinant comparison.  A genus label alone supplies
none of these data.
\end{enumerate}
\end{openproblem}
